\section{Discussion}
\label{sec:discussion}

\paragraph{Evaluate the policy and interface together.}
Strict F1 measures the complete extraction contract; D exposes valid
evidence in the same outputs. Their gap does not additively measure loss
caused by formatting. Identifier-heavy recovery leaves attribute coverage
unresolved. Character overlap cannot charge character false positives for
unanchorable items and must accompany invalid-item counts and exact
penalized precision. None of these measures certifies legal compliance.

\paragraph{Context benefits depend on the output task.}
All ten completed pairs favor local context on D-F1, but both conditions
already expose the same explicit target. The result does not show that
distant evidence is unnecessary. Attention dilution, instruction following
and output allocation are possible explanations, not established mechanisms.
Paired intervals describe document-resampling variation over this synthetic
collection and one execution per condition, not stochastic rerun or
production uncertainty. Shared generation plans may induce dependence;
semantic near-duplicate clusters remain unaudited. T levels use different
documents, not paired rewrites. Constrained generation, revised ownership
or punctuation-tolerant matching would define a new prospective experiment;
we do not retrofit them to improve scores on this test set.

